\label{sec:conclusion}

The 2030 redistricting cycle will be the most technically demanding in the history
of algorithmic redistricting. Three factors combine to create this difficulty:
(1) an increasing number of states require block-group resolution (43 vs.\ 39 in 2020);
(2) Sun Belt population growth requires reapportionment projections for 8 gaining states;
(3) the DIA statutory seed formula requires pre-registration of algorithm configurations
by July 2030, 9 months before the redistricting window opens.

The core recommendation of this paper is simple: begin preparation now.
The 2029--2031 preparation window is achievable but leaves no slack. Any delay
in starting the block-group adjacency construction, the demographic data update,
or the reapportionment projection analysis will compress the preparation window
in ways that may not be recoverable within the statutory redistricting timeline.

\subsection{Roadmap for Q.1--Q.4}

The Q-series papers provide detailed analysis of each preparation task:

\begin{itemize}
\item \textbf{Q.1 --- Data Update Pipeline}: The LODES/ACS update process, including
  interim ACS validation and commuting pattern change analysis from 2021 to 2028.
  Expected magnitude of demographic changes in Sun Belt states.

\item \textbf{Q.2 --- Algorithm Pre-Registration}: State-by-state configuration audit
  for the 2030 pre-registration. Configuration recommendations based on 2020 performance
  and 2030 projection analysis. DIA compliance documentation requirements.

\item \textbf{Q.3 --- Data Infrastructure for 2030}: Detailed analysis of the block-group
  adjacency construction challenge. F.3 resolution rule applied to projected 2030 seat counts.
  43-state block-group adjacency build plan with resource estimates.

\item \textbf{Q.4 --- Reapportionment Projections}: Full Huntington-Hill apportionment
  on Census Bureau medium-series 2030 projections. 8 gainers, 10 losers identified.
  TX prime-factor bisection tree for $k=41$. Cross-census stability analysis.
\end{itemize}

The Q-series is intended to be read by redistricting practitioners, state legislative
staff, and courts that may be asked to evaluate the adequacy of a state's redistricting
preparation. The central message of the series is that algorithmic redistricting
requires infrastructure investment that precedes the redistricting cycle by years,
not months---and that this investment, if made in time, enables redistricting that
is more transparent, more reproducible, and more legally defensible than any
human-drawn alternative.

\bigskip
\noindent\textit{The author thanks the Census Bureau's Redistricting Data Program staff
for public documentation of the P.L.\ 94-171 release timeline and the Census Bureau
medium-series population projections that underlie the Q.4 reapportionment analysis.
All errors are the author's own.}
